% Part: applied-modal-logics
% Chapter: epistemic-logic
% Section: truth-at-w

\documentclass[../../../include/open-logic-section]{subfiles}

\begin{document}

\olfileid{aml}{el}{trw}

\olsection{एखाद्या जगात सत्य असणे}

सामान्य मोडल तर्कशास्त्राप्रमाणेच प्रत्येक ज्ञानविषयक
प्रतिमान त्यात कोणती !!{formula}s कोणत्या जगांत सत्य
मानायची हे ठरवते. संबंधी चिन्हार्थमीमांसेच्या मूलभूत
संकल्पनेसाठी ``प्रतिमान~$\mModel{M}$ मध्ये
!!{formula}~$!A$ हे जग~$w$ येथे सत्य आहे'' हीच संज्ञा
वापरतो. या संबंधाची आगमनाने व्याख्या केली जाते;
मोडल नसलेल्या सर्व संकारकांसाठी ती सामान्य मोडल
प्रकरणातील व्याख्येसारखीच असते.

\begin{defn}\ollabel{defn:mmodels}
  प्रतिमान~$\mModel M$ मध्ये \emph{!!a{formula}~$!A$ हे~$w$ येथे
  सत्य असणे}, म्हणजे $\mSat{M}{!A}[w]$, पुढीलप्रमाणे
  आगमनाने परिभाषित करतो:
  \begin{enumerate}
  \tagitem{prvFalse}{\indcase{!A}{\lfalse}{$\mSat{M}{\lfalse}[w]$ कधीच नसते}.}{}
  \tagitem{prvTrue}{\indcase{!A}{\ltrue}{$\mSat{M}{\ltrue}[w]$ नेहमी असते}.}{}
  \item $\mSat{M}{p}[w]$ तेव्हा आणि केवळ तेव्हाच, जेव्हा $w \in V(p)$
  \tagitem{prvNot}{\indcase{!A}{\lnot !B}{$\mSat{M}{\indfrm}[w]$
    तेव्हा आणि केवळ तेव्हाच, जेव्हा $\mSat/{M}{!B}[w]$}.}{}
  \tagitem{prvAnd}{\indcase{!A}{(!B \land !C)}{$\mSat{M}{\indfrm}[w]$
    तेव्हा आणि केवळ तेव्हाच, जेव्हा $\mSat{M}{!B}[w]$ आणि
    $\mSat{M}{!C}[w]$}.}{}
  \tagitem{prvOr}{\indcase{!A}{(!B \lor !C)}{$\mSat{M}{\indfrm}[w]$
    तेव्हा आणि केवळ तेव्हाच, जेव्हा $\mSat{M}{!B}[w]$ किंवा
    $\mSat{M}{!C}[w]$} (किंवा दोन्ही).}{}
  \tagitem{prvIf}{\indcase{!A}{(!B \lif !C)}{$\mSat{M}{\indfrm}[w]$
    तेव्हा आणि केवळ तेव्हाच, जेव्हा $\mSat/{M}{!B}[w]$ किंवा
    $\mSat{M}{!C}[w]$}.}{}
  \tagitem{prvIff}{\indcase{!A}{(!B \liff !C)}{$\mSat{M}{\indfrm}[w]$
    तेव्हा आणि केवळ तेव्हाच, जेव्हा $\mSat{M}{!B}[w]$ आणि
    $\mSat{M}{!C}[w]$ ही दोन्ही असतात, किंवा
    $\mSat{M}{!B}[w]$ आणि $\mSat{M}{!C}[w]$ यांपैकी
    कोणतेही नसते}.}{}
  \item\ollabel{defn:sub:mmodels-box}
    \indcase{!A}{\Knows_a !B}{$\mSat{M}{\indfrm}[w]$ तेव्हा आणि
    केवळ तेव्हाच, जेव्हा $R_a ww'$ असलेल्या प्रत्येक $w' \in W$
    साठी $\mSat{M}{!B}[w']$}
  \end{enumerate}
\end{defn}

इथे मात्र प्राप्यता संबंधांवर घालायच्या अटींचा विचार
करावा लागतो. खंड~\olref{defn:sub:mmodels-box} नुसार,
$R_a ww'$ असलेले एकही~$w'$ नसेल, तेव्हा
!!a{formula}~$\Knows_a !B$ हे~$w$ येथे सत्य असते.
हा खंड सामान्य मोडल तर्कशास्त्रातील खंडासारखाच
आहे: एखाद्या जगाला उत्तरवर्ती जगे नसतील, तर सर्व
$\Box$-सूत्रे तेथे रिक्तपणे सत्य असतात. $\Knows$ हा
\emph{ज्ञान} दर्शवतो असे मानल्यास हे अत्यंत
प्रतिसहजभासी वाटते. कारण एखाद्या कर्त्याला एकाच
वेळी $!A$ आणि $\lnot !A$ ही दोन्ही माहीत असू शकतील,
अशी \emph{कोणतीही} परिस्थिती नसते, असे आपण
सहसा मानतो.

एक उपाय म्हणजे ज्ञानविषयक तर्कशास्त्रातील प्राप्यता
संबंध नेहमी \emph{परावर्ती} असेल याची खात्री करणे.
यातून वास्तविक जग कर्त्याच्या माहितीशी सुसंगत आहे,
ही कल्पना ढोबळमानाने व्यक्त होते. वस्तुतः ज्ञानविषयक
तर्कशास्त्रांत सहसा S5 वापरतात; पण $\Knows_a$ संबंधाने
नेमके काय दर्शवायचे आहे यानुसार काही जण त्याहून
दुर्बल व्यवस्था वापरू शकतात.

\begin{figure}
  \begin{center}
    \begin{tikzpicture}[modal]
      \node[world] (w1) [label={[align=right]right:\mTrue{p}\\ \mFalse{q}}]{$w_1$};
      \node[world] (w2) [label={[align=right]right:\mTrue{p}\\ \mTrue{q}},
        above right=of w1]{$w_2$};
      \node[world] (w3) [label={[align=right]right:\mFalse{p}\\ \mFalse{q}},
        below right=of w1] {$w_3$};
      \draw[<->] (w1) to node {$a$} (w2);
      \draw[->,loop left] (w1) to node {$a,b$} (w1);
      \draw[<->] (w1) to [swap] node {$b$} (w3);
      \draw[->,loop above] (w2) to node {$a,b$} (w2) ;
      \draw[->,loop below] (w3) to node {$a,b$} (w3) ;
    \end{tikzpicture}
  \end{center}
  \caption{साधे ज्ञानविषयक प्रतिमान.}
  \ollabel{fig:simple}
\end{figure}

\begin{prob}
  \olref[aml][el][trw]{fig:simple} मधील प्रतिमानात
  पुढीलपैकी कोणती विधाने सत्य आहेत ते ठरवा:
  \begin{enumerate}
    \item $\mSat{M}{\lnot q}[w_1]$;
    \item $\mSat{M}{\Knows_a \lnot q}[w_1]$;
    \item $\mSat{M}{\Knows_b \lnot q}[w_1]$;
    \item $\mSat{M}{\Knows_b q \lor \Knows_b \lnot q}[w_2]$;
    \item $\mSat{M}{\Knows_a( \Knows_b q \lor \Knows_b \lnot q)}[w_2]$;
    \item $\mSat{M}{\EKnows_{\{a,b\}} \lnot q}[w_3]$;
    \end{enumerate}
\end{prob}

जगात सत्य असण्याची मूलभूत व्याख्या दिल्यानंतर,
मोडल वैधता आणि तार्किक निष्पन्नता यांसारख्या सामान्य
मोडल तर्कशास्त्रातील इतर चिन्हार्थमीमांसक संकल्पनाही
इथे लागू होतात; फक्त मोडल संकारकांच्या अर्थलक्षणाकडे
पाहण्याचा हा नवा दृष्टिकोन वापरावा लागतो.

आता सामायिक-ज्ञान संकारक~$\CKnows_G$ साठी सत्यतेच्या
अटीही देता येतात. \olref[sfr][rel][ops]{sec} मधील
व्याख्येनुसार, संबंध~$R$ चे \emph{संक्रमक संवरण}~$R^+$
पुढीलप्रमाणे असते:
\begin{align*}
  R^+ &= \bigcup_{ n \in \mathbb{N}} R^n,
\intertext{येथे}
  R^0 & = R \text{ आणि}\\
  R^{n+1} & = \Setabs{\tuple{x, z}}{\exists y (R^n xy \land Ryz)}.
\end{align*}
मग $G$ हा कर्त्यांचा समूह असेल, तर सर्व कर्त्यांच्या
प्राप्यता संबंधांच्या संघाचे संक्रमक संवरण म्हणून
$R_G = ( \bigcup_{b
\in G} R_b )^+$ परिभाषित करतो.

\begin{defn}
$G' \subseteq G$ असेल, तर $\mSat{M}{\CKnows_{G'} !A}[ w]$
तेव्हा आणि केवळ तेव्हाच, जेव्हा $R_{G'} w w'$ असलेल्या
प्रत्येक $w'$ साठी $\mSat{M}{!A}[w']$.
\end{defn}

\end{document}
